% =====================
% Chapter 16: Data Visualization with Pandas
% =====================
% Requires: \h, \hh, \hhh, \code macros; 4 spaces per indent

\h{Data Visualization with Pandas and Matplotlib}

Imagine presenting your building analysis to a client. You've calculated that 40\% of the space is offices, 35\% is public areas, and 25\% is service spaces. While Pandas can create quick charts with `.plot()`, combining it with matplotlib gives you complete control over every aspect of your visualization---from colors and fonts to annotations and custom layouts. This is the difference between a quick sketch and a presentation-ready drawing.

\emph{Matplotlib} is Python's foundational plotting library, and Pandas is built on top of it. When you use Pandas' `.plot()` method, matplotlib is working behind the scenes. By learning both together, you get the best of both worlds: Pandas' convenience for quick plots and matplotlib's power for customization.

\hh{Real-world architectural applications}

Professional architectural visualizations require precise control over appearance and layout:

\begin{itemize}
\item \textbf{Client presentations:} Branded colors, custom fonts, professional annotations
\item \textbf{Technical reports:} Multiple subplots, detailed labels, reference lines
\item \textbf{Competition boards:} High-resolution exports, consistent styling, complex layouts
\item \textbf{BIM integration:} Automated chart generation with corporate templates
\item \textbf{Publication figures:} Journal-ready formatting, precise dimensions, vector exports
\item \textbf{Dashboard displays:} Real-time updates, interactive elements, custom layouts
\end{itemize}

\begin{tcolorbox}[title=Focus of this chapter:,colframe=orange,colback=white,arc=2mm]
This chapter teaches you to create professional visualizations combining Pandas and matplotlib:
\begin{itemize}
\item \textbf{Setup:} Understanding the matplotlib + Pandas relationship
\item \textbf{Basic plotting:} Using both Pandas shortcuts and matplotlib control
\item \textbf{Customization:} Fonts, colors, sizes, and styles
\item \textbf{Annotations:} Adding text, arrows, and highlights
\item \textbf{Subplots:} Creating complex multi-chart layouts
\item \textbf{Professional styling:} Publication-ready charts
\item \textbf{Export:} High-quality saves for reports and presentations
\end{itemize}
\end{tcolorbox}

% ====================================================================
% PART I – UNDERSTANDING MATPLOTLIB + PANDAS
% ====================================================================

\hh{Setting up matplotlib with Pandas}

Let's understand how Pandas and matplotlib work together, starting with the basics.

\code{c16_setup}{python}{
    import pandas as pd
    import matplotlib.pyplot as plt
    import numpy as np
    
    # Set up matplotlib for better display
    plt.style.use('default')  # Can try: 'seaborn', 'ggplot', 'bmh'
    
    # For high-DPI displays (optional)
    plt.rcParams['figure.dpi'] = 100
    plt.rcParams['savefig.dpi'] = 300
    
    # Create sample data
    room_areas = pd.Series([45, 32, 28, 38, 25],
                           index=["Lobby", "Office", "Studio", "Gallery", "Storage"])
    
    # Method 1: Pure Pandas (matplotlib works behind scenes)
    room_areas.plot(kind='bar', title='Room Areas')
    plt.ylabel('Area (m^2)')  # Adding matplotlib customization
    plt.show()
    
    # Method 2: Get the axes object for more control
    fig, ax = plt.subplots(figsize=(8, 5))
    room_areas.plot(kind='bar', ax=ax, color='steelblue')
    ax.set_title('Room Areas in Building A', fontsize=14, fontweight='bold')
    ax.set_ylabel('Area (m^2)', fontsize=12)
    ax.set_xlabel('Room Type', fontsize=12)
    ax.grid(True, alpha=0.3)
    plt.tight_layout()
    plt.show()
    
    print("Two approaches:")
    print("1. Quick plot with Pandas + minor matplotlib tweaks")
    print("2. Full control with matplotlib figure and axes")
    
    # Outcome:
    # Two approaches:
    # 1. Quick plot with Pandas + minor matplotlib tweaks
    # 2. Full control with matplotlib figure and axes
    # [Two bar charts appear - second one more polished]
}{Understanding the Pandas-matplotlib relationship}

\textbf{Key concepts:}
\begin{itemize}
\item \textbf{Figure:} The entire image/window containing plots
\item \textbf{Axes:} The actual plot area within the figure
\item \textbf{plt:} Module-level functions for quick plotting
\item Pandas \c{.plot()} returns an axes object you can customize
\item \c{plt.subplots()} creates figure and axes for full control
\end{itemize}

\hh{The matplotlib object hierarchy}

Understanding matplotlib's structure helps you customize effectively.

\code{c16_hierarchy}{python}{
    import pandas as pd
    import matplotlib.pyplot as plt
    
    # Create a figure with labeled components
    fig, ax = plt.subplots(figsize=(10, 6))
    
    # Add a simple plot
    data = pd.Series([30, 45, 25, 40], index=['Q1', 'Q2', 'Q3', 'Q4'])
    data.plot(kind='line', ax=ax, marker='o', linewidth=2)
    
    # Customize different components
    # Figure level
    fig.suptitle('Building Energy Consumption', fontsize=16, y=0.98)
    
    # Axes level
    ax.set_title('Quarterly Breakdown 2024', fontsize=12, pad=20)
    ax.set_xlabel('Quarter', fontsize=11, fontweight='bold')
    ax.set_ylabel('Energy (MWh)', fontsize=11, fontweight='bold')
    ax.set_ylim(0, 50)
    
    # Grid and spines
    ax.grid(True, linestyle='--', alpha=0.7)
    ax.spines['top'].set_visible(False)
    ax.spines['right'].set_visible(False)
    
    # Add annotations
    ax.annotate('Peak Usage', xy=('Q2', 45), xytext=('Q3', 47),
                arrowprops=dict(arrowstyle='->', color='red'),
                fontsize=10, color='red')
    
    # Add reference line
    ax.axhline(y=35, color='green', linestyle=':', label='Target')
    ax.legend()
    
    plt.tight_layout()
    plt.show()
    
    print("Matplotlib components customized:")
    print("- Figure: Overall container and main title")
    print("- Axes: Plot area, labels, limits, grid")
    print("- Spines: Border lines around plot")
    print("- Annotations: Text and arrows")
    print("- Legend: Automatically from labels")
    
    # Outcome:
    # Matplotlib components customized:
    # - Figure: Overall container and main title
    # - Axes: Plot area, labels, limits, grid
    # - Spines: Border lines around plot
    # - Annotations: Text and arrows
    # - Legend: Automatically from labels
    # [Polished line chart with annotations appears]
}{Understanding matplotlib's component hierarchy}

% ====================================================================
% PART II – ENHANCED BAR CHARTS
% ====================================================================

\hh{Professional bar charts with full customization}

Let's create publication-quality bar charts by combining Pandas data handling with matplotlib styling.

\code{c16_bar_enhanced}{python}{
    import pandas as pd
    import matplotlib.pyplot as plt
    import numpy as np
    
    # Building floor areas data
    df = pd.DataFrame({
        'Floor': ['Basement', 'Ground', 'First', 'Second', 'Third', 'Roof'],
        'Total_Area': [450, 580, 560, 540, 480, 120],
        'Usable_Area': [380, 520, 510, 495, 430, 80]
    })
    df.set_index('Floor', inplace=True)
    
    # Create figure with custom size
    fig, (ax1, ax2) = plt.subplots(1, 2, figsize=(14, 6))
    
    # Left plot: Simple bars with custom colors
    colors = ['#2E86AB', '#A23B72', '#F18F01', '#C73E1D', '#6A994E', '#BC4B51']
    df['Total_Area'].plot(kind='bar', ax=ax1, color=colors, width=0.7)
    
    ax1.set_title('Total Area by Floor', fontsize=14, fontweight='bold', pad=20)
    ax1.set_ylabel('Area (m^2)', fontsize=12)
    ax1.set_xlabel('')
    ax1.set_xticklabels(ax1.get_xticklabels(), rotation=45, ha='right')
    
    # Add value labels on bars
    for i, (idx, value) in enumerate(df['Total_Area'].items()):
        ax1.text(i, value + 10, f'{value:.0f}', ha='center', va='bottom',
                fontweight='bold', fontsize=10)
    
    # Add average line
    avg = df['Total_Area'].mean()
    ax1.axhline(y=avg, color='red', linestyle='--', alpha=0.7, 
               label=f'Average: {avg:.0f} m^2')
    ax1.legend()
    
    # Right plot: Grouped bars with custom styling
    x = np.arange(len(df.index))
    width = 0.35
    
    bars1 = ax2.bar(x - width/2, df['Total_Area'], width, 
                    label='Total', color='#2E86AB', alpha=0.8)
    bars2 = ax2.bar(x + width/2, df['Usable_Area'], width,
                    label='Usable', color='#F18F01', alpha=0.8)
    
    ax2.set_title('Total vs Usable Area', fontsize=14, fontweight='bold', pad=20)
    ax2.set_ylabel('Area (m^2)', fontsize=12)
    ax2.set_xticks(x)
    ax2.set_xticklabels(df.index, rotation=45, ha='right')
    ax2.legend(loc='upper right')
    ax2.grid(True, axis='y', alpha=0.3)
    
    # Add efficiency percentages
    for i, (total, usable) in enumerate(zip(df['Total_Area'], df['Usable_Area'])):
        efficiency = (usable/total) * 100
        ax2.text(i, max(total, usable) + 20, f'{efficiency:.0f}%',
                ha='center', fontsize=9, style='italic')
    
    plt.suptitle('Building Floor Analysis', fontsize=16, fontweight='bold', y=1.02)
    plt.tight_layout()
    plt.show()
    
    print("Professional bar charts created with:")
    print("- Custom colors and transparency")
    print("- Value labels on bars")
    print("- Reference lines and averages")
    print("- Efficiency percentages")
    print("- Professional typography")
    
    # Outcome:
    # Professional bar charts created with:
    # - Custom colors and transparency
    # - Value labels on bars
    # - Reference lines and averages
    # - Efficiency percentages
    # - Professional typography
    # [Two professional bar charts side by side]
}{Creating publication-quality bar charts}

\hhh{Stacked bars with patterns and textures}

For black-and-white printing or accessibility, patterns can replace colors.

\code{c16_stacked_patterns}{python}{
    import pandas as pd
    import matplotlib.pyplot as plt
    import numpy as np
    
    # Space allocation data
    df = pd.DataFrame({
        'Floor': ['Ground', 'First', 'Second', 'Third'],
        'Office': [200, 250, 180, 160],
        'Meeting': [80, 60, 90, 70],
        'Public': [150, 100, 120, 140],
        'Service': [50, 40, 45, 55]
    })
    df.set_index('Floor', inplace=True)
    
    fig, (ax1, ax2) = plt.subplots(1, 2, figsize=(14, 6))
    
    # Left: Color version
    df.plot(kind='bar', stacked=True, ax=ax1, 
            colormap='Set3', width=0.7)
    ax1.set_title('Space Allocation by Floor (Color)', fontsize=13, fontweight='bold')
    ax1.set_ylabel('Area (m^2)', fontsize=11)
    ax1.set_xlabel('Floor', fontsize=11)
    ax1.legend(title='Space Type', bbox_to_anchor=(1.05, 1), loc='upper left')
    ax1.set_xticklabels(ax1.get_xticklabels(), rotation=0)
    
    # Right: Pattern version (for B&W printing)
    patterns = ['///', '\\\\\\', '|||', '---']
    colors = ['white', 'lightgray', 'darkgray', 'black']
    
    bottom = np.zeros(len(df.index))
    for i, col in enumerate(df.columns):
        ax2.bar(df.index, df[col], bottom=bottom, 
               label=col, color=colors[i], 
               edgecolor='black', linewidth=1,
               hatch=patterns[i])
        bottom += df[col]
    
    ax2.set_title('Space Allocation by Floor (B&W)', fontsize=13, fontweight='bold')
    ax2.set_ylabel('Area (m^2)', fontsize=11)
    ax2.set_xlabel('Floor', fontsize=11)
    ax2.legend(title='Space Type', bbox_to_anchor=(1.05, 1), loc='upper left')
    
    # Add total labels on top
    totals = df.sum(axis=1)
    for i, (floor, total) in enumerate(totals.items()):
        ax1.text(i, total + 10, f'{total:.0f} m^2', 
                ha='center', fontweight='bold')
        ax2.text(i, total + 10, f'{total:.0f} m^2', 
                ha='center', fontweight='bold')
    
    plt.suptitle('Stacked Bar Comparison: Color vs B&W', 
                fontsize=14, fontweight='bold', y=1.02)
    plt.tight_layout()
    plt.show()
    
    print("Stacked bars with:")
    print("- Color version for digital display")
    print("- Pattern version for B&W printing")
    print("- Total values on top")
    print("- Clear legends and labels")
    
    # Outcome:
    # Stacked bars with:
    # - Color version for digital display
    # - Pattern version for B&W printing
    # - Total values on top
    # - Clear legends and labels
    # [Two stacked bar charts - color and B&W versions]
}{Creating accessible visualizations with patterns}

% ====================================================================
% PART III – ADVANCED LINE CHARTS AND TIME SERIES
% ====================================================================

\hh{Time series with multiple scales and annotations}

Line charts become powerful when showing trends with proper annotations and multiple data series.

\code{c16_timeseries_advanced}{python}{
    import pandas as pd
    import matplotlib.pyplot as plt
    import numpy as np
    
    # Create monthly building performance data
    months = pd.date_range('2024-01', '2024-12', freq='M')
    df = pd.DataFrame({
        'Month': months,
        'Energy_MWh': [45, 42, 38, 32, 28, 25, 26, 27, 30, 35, 40, 44],
        'Cost_kEUR': [5.8, 5.4, 4.9, 4.1, 3.6, 3.2, 3.3, 3.5, 3.9, 4.5, 5.2, 5.7],
        'Occupancy': [145, 152, 158, 165, 170, 162, 155, 148, 168, 172, 160, 150],
        'Comfort_Score': [72, 74, 78, 82, 85, 88, 86, 87, 83, 79, 75, 71]
    })
    df.set_index('Month', inplace=True)
    
    # Create figure with specific layout
    fig = plt.figure(figsize=(15, 10))
    
    # Main plot: Energy and Cost with dual axes
    ax1 = plt.subplot(2, 1, 1)
    
    # Primary y-axis: Energy
    color1 = '#2E86AB'
    ax1.plot(df.index, df['Energy_MWh'], marker='o', linewidth=2.5,
            color=color1, label='Energy Consumption')
    ax1.set_xlabel('Month', fontsize=11)
    ax1.set_ylabel('Energy (MWh)', color=color1, fontsize=11)
    ax1.tick_params(axis='y', labelcolor=color1)
    ax1.grid(True, alpha=0.3)
    
    # Secondary y-axis: Cost
    ax2 = ax1.twinx()
    color2 = '#A23B72'
    ax2.plot(df.index, df['Cost_kEUR'], marker='s', linewidth=2.5,
            color=color2, linestyle='--', label='Cost')
    ax2.set_ylabel('Cost (kEUR)', color=color2, fontsize=11)
    ax2.tick_params(axis='y', labelcolor=color2)
    
    # Add title and legends
    ax1.set_title('Building Performance Metrics 2024', 
                 fontsize=14, fontweight='bold', pad=20)
    
    # Combine legends
    lines1, labels1 = ax1.get_legend_handles_labels()
    lines2, labels2 = ax2.get_legend_handles_labels()
    ax1.legend(lines1 + lines2, labels1 + labels2, loc='upper right')
    
    # Add annotations for significant events
    ax1.annotate('Summer\nMinimum', 
                xy=(df.index[5], df['Energy_MWh'].iloc[5]),
                xytext=(df.index[3], 20),
                arrowprops=dict(arrowstyle='->', color='red', lw=1.5),
                fontsize=10, color='red', ha='center')
    
    ax1.annotate('Winter\nPeak', 
                xy=(df.index[0], df['Energy_MWh'].iloc[0]),
                xytext=(df.index[2], 48),
                arrowprops=dict(arrowstyle='->', color='red', lw=1.5),
                fontsize=10, color='red', ha='center')
    
    # Shade seasons
    ax1.axvspan(df.index[5], df.index[7], alpha=0.2, color='yellow', label='Summer')
    ax1.axvspan(df.index[11], df.index[11], alpha=0.2, color='lightblue')
    ax1.axvspan(df.index[0], df.index[1], alpha=0.2, color='lightblue', label='Winter')
    
    # Bottom plot: Occupancy and Comfort
    ax3 = plt.subplot(2, 1, 2)
    
    # Create filled area chart
    ax3.fill_between(df.index, 0, df['Occupancy'], 
                    alpha=0.3, color='green', label='Occupancy')
    ax3.plot(df.index, df['Occupancy'], color='green', linewidth=2)
    
    # Add comfort score on secondary axis
    ax4 = ax3.twinx()
    ax4.plot(df.index, df['Comfort_Score'], color='orange', 
            marker='D', linewidth=2, label='Comfort Score')
    
    ax3.set_xlabel('Month', fontsize=11)
    ax3.set_ylabel('Occupancy (people)', color='green', fontsize=11)
    ax4.set_ylabel('Comfort Score (%)', color='orange', fontsize=11)
    ax3.tick_params(axis='y', labelcolor='green')
    ax4.tick_params(axis='y', labelcolor='orange')
    ax3.grid(True, alpha=0.3)
    
    ax3.set_title('Occupancy and Comfort Levels', 
                 fontsize=12, fontweight='bold', pad=15)
    
    # Format x-axis
    import matplotlib.dates as mdates
    ax1.xaxis.set_major_formatter(mdates.DateFormatter('%b'))
    ax3.xaxis.set_major_formatter(mdates.DateFormatter('%b'))
    
    plt.tight_layout()
    plt.show()
    
    print("Advanced time series features:")
    print("- Dual y-axes for different scales")
    print("- Seasonal shading")
    print("- Event annotations")
    print("- Filled area charts")
    print("- Multiple related plots")
    
    # Outcome:
    # Advanced time series features:
    # - Dual y-axes for different scales
    # - Seasonal shading
    # - Event annotations
    # - Filled area charts
    # - Multiple related plots
    # [Complex multi-panel time series visualization]
}{Creating professional time series visualizations}

% ====================================================================
% PART IV – STATISTICAL PLOTS WITH STYLE
% ====================================================================

\hh{Enhanced distribution plots}

Combining Pandas statistical plots with matplotlib customization for publication-ready figures.

\code{c16_distributions_styled}{python}{
    import pandas as pd
    import matplotlib.pyplot as plt
    import numpy as np
    
    # Generate realistic room size data
    np.random.seed(42)
    office_sizes = np.random.normal(25, 5, 100)
    meeting_sizes = np.random.normal(35, 8, 80)
    public_sizes = np.random.normal(60, 15, 60)
    
    # Create figure with subplots
    fig, axes = plt.subplots(2, 2, figsize=(14, 10))
    fig.suptitle('Room Size Distribution Analysis', 
                fontsize=16, fontweight='bold', y=1.02)
    
    # 1. Overlapping histograms with transparency
    ax1 = axes[0, 0]
    ax1.hist(office_sizes, bins=20, alpha=0.5, label='Office', 
            color='blue', edgecolor='black')
    ax1.hist(meeting_sizes, bins=20, alpha=0.5, label='Meeting', 
            color='green', edgecolor='black')
    ax1.hist(public_sizes, bins=20, alpha=0.5, label='Public', 
            color='red', edgecolor='black')
    
    ax1.set_title('Distribution Comparison', fontsize=12, fontweight='bold')
    ax1.set_xlabel('Room Size (m2)', fontsize=10)
    ax1.set_ylabel('Frequency', fontsize=10)
    ax1.legend()
    ax1.grid(True, alpha=0.3)
    
    # Add statistics text
    stats_text = f'Office: mu={office_sizes.mean():.1f}, sigma={office_sizes.std():.1f}\n'
    stats_text += f'Meeting: mu={meeting_sizes.mean():.1f}, sigma={meeting_sizes.std():.1f}\n'
    stats_text += f'Public: mu={public_sizes.mean():.1f}, sigma={public_sizes.std():.1f}'
    ax1.text(0.02, 0.98, stats_text, transform=ax1.transAxes,
            fontsize=9, verticalalignment='top',
            bbox=dict(boxstyle='round', facecolor='wheat', alpha=0.5))
    
    # 2. Box plot with violin plot overlay
    ax2 = axes[0, 1]
    df_sizes = pd.DataFrame({
        'Office': office_sizes,
        'Meeting': np.append(meeting_sizes, [np.nan]*20),  # Pad with NaN
        'Public': np.append(public_sizes, [np.nan]*40)
    })
    
    # Create box plot
    bp = ax2.boxplot([office_sizes, meeting_sizes, public_sizes],
                     labels=['Office', 'Meeting', 'Public'],
                     patch_artist=True, widths=0.5)
    
    # Color the boxes
    colors = ['lightblue', 'lightgreen', 'lightcoral']
    for patch, color in zip(bp['boxes'], colors):
        patch.set_facecolor(color)
        patch.set_alpha(0.7)
    
    # Customize whiskers and medians
    for whisker in bp['whiskers']:
        whisker.set(color='gray', linewidth=1.5, linestyle=':')
    for median in bp['medians']:
        median.set(color='red', linewidth=2)
    
    ax2.set_title('Statistical Summary', fontsize=12, fontweight='bold')
    ax2.set_ylabel('Room Size (m2)', fontsize=10)
    ax2.grid(True, axis='y', alpha=0.3)
    
    # 3. Scatter plot with trend line
    ax3 = axes[1, 0]
    
    # Create correlated data
    areas = np.random.uniform(20, 100, 150)
    costs = areas * 500 + np.random.normal(0, 5000, 150)
    
    # Scatter plot
    scatter = ax3.scatter(areas, costs, c=areas, cmap='viridis', 
                         alpha=0.6, s=50, edgecolors='black', linewidth=0.5)
    
    # Add trend line
    z = np.polyfit(areas, costs, 1)
    p = np.poly1d(z)
    ax3.plot(np.sort(areas), p(np.sort(areas)), "r--", alpha=0.8, 
            label=f'y = {z[0]:.1f}x + {z[1]:.0f}')
    
    # Add colorbar
    cbar = plt.colorbar(scatter, ax=ax3)
    cbar.set_label('Area (m2)', fontsize=10)
    
    ax3.set_title('Area vs Cost Relationship', fontsize=12, fontweight='bold')
    ax3.set_xlabel('Area (m2)', fontsize=10)
    ax3.set_ylabel('Cost (EUR)', fontsize=10)
    ax3.legend()
    ax3.grid(True, alpha=0.3)
    
    # Calculate and display R-squared
    correlation = np.corrcoef(areas, costs)[0, 1]
    r_squared = correlation ** 2
    ax3.text(0.05, 0.95, f'R-squared = {r_squared:.3f}', transform=ax3.transAxes,
            fontsize=10, verticalalignment='top',
            bbox=dict(boxstyle='round', facecolor='yellow', alpha=0.5))
    
    # 4. Pie chart with exploded slices
    ax4 = axes[1, 1]
    
    # Calculate space allocation
    space_types = ['Office', 'Meeting', 'Public', 'Service', 'Circulation']
    sizes = [35, 20, 25, 10, 10]
    explode = (0.1, 0, 0, 0, 0)  # Explode office slice
    colors_pie = ['#FF9999', '#66B2FF', '#99FF99', '#FFCC99', '#FF99CC']
    
    wedges, texts, autotexts = ax4.pie(sizes, explode=explode, labels=space_types,
                                        colors=colors_pie, autopct='%1.1f%%',
                                        shadow=True, startangle=90)
    
    # Customize text
    for text in texts:
        text.set_fontsize(10)
    for autotext in autotexts:
        autotext.set_color('white')
        autotext.set_fontweight('bold')
        autotext.set_fontsize(9)
    
    ax4.set_title('Space Allocation', fontsize=12, fontweight='bold')
    
    # Add total area text
    total_area = 2500  # m2
    ax4.text(0, -1.3, f'Total Area: {total_area:,} m2', 
            ha='center', fontsize=10, fontweight='bold')
    
    plt.tight_layout()
    plt.show()
    
    print("Statistical visualization dashboard created with:")
    print("- Overlapping histograms with statistics")
    print("- Customized box plots")
    print("- Scatter plot with trend line and R-squared")
    print("- Enhanced pie chart")
    
    # Outcome:
    # Statistical visualization dashboard created with:
    # - Overlapping histograms with statistics
    # - Customized box plots
    # - Scatter plot with trend line and R-squared
    # - Enhanced pie chart
    # [Four-panel statistical dashboard]
}{Creating a statistical visualization dashboard}

% ====================================================================
% PART V – COMPLEX LAYOUTS AND SUBPLOTS
% ====================================================================

\hh{Creating complex dashboard layouts}

Professional reports often need multiple related visualizations arranged precisely.

\code{c16_complex_layout}{python}{
    import pandas as pd
    import matplotlib.pyplot as plt
    import matplotlib.gridspec as gridspec
    import numpy as np
    
    # Create comprehensive building data
    np.random.seed(42)
    
    # Monthly data
    months = ['Jan', 'Feb', 'Mar', 'Apr', 'May', 'Jun',
              'Jul', 'Aug', 'Sep', 'Oct', 'Nov', 'Dec']
    
    energy_data = pd.DataFrame({
        'Heating': [450, 400, 300, 150, 50, 0, 0, 0, 30, 180, 350, 420],
        'Cooling': [0, 0, 20, 80, 150, 280, 350, 320, 180, 50, 0, 0],
        'Lighting': [180, 170, 150, 130, 120, 110, 110, 115, 135, 160, 175, 185],
        'Equipment': [200, 200, 200, 200, 200, 200, 200, 200, 200, 200, 200, 200]
    }, index=months)
    
    # Create figure with GridSpec for complex layout
    fig = plt.figure(figsize=(16, 10))
    gs = gridspec.GridSpec(3, 3, figure=fig, hspace=0.3, wspace=0.3)
    
    # Main title
    fig.suptitle('Building Performance Dashboard - 2024', 
                fontsize=18, fontweight='bold', y=0.98)
    
    # 1. Large plot: Stacked area chart (spans 2 columns)
    ax1 = fig.add_subplot(gs[0, :2])
    energy_data.plot(kind='area', stacked=True, ax=ax1, 
                     alpha=0.7, colormap='tab20c')
    ax1.set_title('Monthly Energy Consumption by Type', 
                 fontsize=12, fontweight='bold')
    ax1.set_xlabel('Month')
    ax1.set_ylabel('Energy (kWh)')
    ax1.legend(loc='upper left', ncol=4, framealpha=0.9)
    ax1.grid(True, alpha=0.3)
    
    # Add total line
    total = energy_data.sum(axis=1)
    ax1.plot(total.index, total.values, 'k-', linewidth=2, 
            label='Total', linestyle='--')
    
    # 2. Pie chart: Annual distribution
    ax2 = fig.add_subplot(gs[0, 2])
    annual_totals = energy_data.sum()
    colors = plt.cm.Set3(np.arange(len(annual_totals)))
    wedges, texts, autotexts = ax2.pie(annual_totals, 
                                        labels=annual_totals.index,
                                        autopct='%1.1f%%',
                                        colors=colors,
                                        startangle=90)
    ax2.set_title('Annual Energy\nDistribution', 
                 fontsize=11, fontweight='bold')
    
    # Make percentage text smaller
    for autotext in autotexts:
        autotext.set_fontsize(8)
    
    # 3. Bar chart: Quarterly comparison
    ax3 = fig.add_subplot(gs[1, 0])
    quarterly = pd.DataFrame({
        'Q1': energy_data.iloc[:3].sum(),
        'Q2': energy_data.iloc[3:6].sum(),
        'Q3': energy_data.iloc[6:9].sum(),
        'Q4': energy_data.iloc[9:].sum()
    }).T
    
    quarterly.plot(kind='bar', ax=ax3, width=0.8)
    ax3.set_title('Quarterly Comparison', fontsize=11, fontweight='bold')
    ax3.set_xlabel('Quarter')
    ax3.set_ylabel('Energy (kWh)')
    ax3.set_xticklabels(ax3.get_xticklabels(), rotation=0)
    ax3.legend(loc='upper right', fontsize=8)
    ax3.grid(True, axis='y', alpha=0.3)
    
    # 4. Line chart with markers: Efficiency metrics
    ax4 = fig.add_subplot(gs[1, 1])
    efficiency = pd.Series([82, 84, 86, 88, 90, 91, 89, 88, 87, 85, 83, 81],
                          index=months)
    efficiency.plot(ax=ax4, marker='o', linewidth=2, color='green')
    ax4.set_title('System Efficiency', fontsize=11, fontweight='bold')
    ax4.set_ylabel('Efficiency (%)')
    ax4.set_ylim(75, 95)
    ax4.grid(True, alpha=0.3)
    ax4.fill_between(efficiency.index, 85, efficiency.values,
                    where=(efficiency >= 85), alpha=0.3, color='green',
                    interpolate=True, label='Above Target')
    ax4.axhline(y=85, color='red', linestyle=':', label='Target: 85%')
    ax4.legend(fontsize=8)
    
    # 5. Heatmap: Daily patterns
    ax5 = fig.add_subplot(gs[1, 2])
    
    # Create sample hourly data (24 hours x 7 days)
    hourly_data = np.random.randn(24, 7) * 10 + 50
    im = ax5.imshow(hourly_data, cmap='YlOrRd', aspect='auto')
    ax5.set_title('Weekly Usage Pattern', fontsize=11, fontweight='bold')
    ax5.set_xlabel('Day of Week')
    ax5.set_ylabel('Hour of Day')
    ax5.set_xticks(range(7))
    ax5.set_xticklabels(['Mon', 'Tue', 'Wed', 'Thu', 'Fri', 'Sat', 'Sun'])
    ax5.set_yticks(range(0, 24, 4))
    ax5.set_yticklabels(['00:00', '04:00', '08:00', '12:00', '16:00', '20:00'])
    
    # Add colorbar
    cbar = plt.colorbar(im, ax=ax5, fraction=0.046, pad=0.04)
    cbar.set_label('kWh', fontsize=9)
    
    # 6. Box plot: Distribution by season
    ax6 = fig.add_subplot(gs[2, :2])
    
    seasonal_data = pd.DataFrame({
        'Winter': energy_data.iloc[[11, 0, 1]].sum(axis=1),
        'Spring': energy_data.iloc[[2, 3, 4]].sum(axis=1),
        'Summer': energy_data.iloc[[5, 6, 7]].sum(axis=1),
        'Fall': energy_data.iloc[[8, 9, 10]].sum(axis=1)
    })
    
    bp = seasonal_data.boxplot(ax=ax6, patch_artist=True,
                               boxprops=dict(facecolor='lightblue'))
    ax6.set_title('Seasonal Energy Distribution', fontsize=11, fontweight='bold')
    ax6.set_ylabel('Total Energy (kWh)')
    ax6.grid(True, axis='y', alpha=0.3)
    
    # 7. Summary statistics table
    ax7 = fig.add_subplot(gs[2, 2])
    ax7.axis('tight')
    ax7.axis('off')
    
    # Calculate summary statistics
    summary_data = [
        ['Metric', 'Value'],
        ['Annual Total', f'{energy_data.sum().sum():,.0f} kWh'],
        ['Monthly Avg', f'{energy_data.sum().sum()/12:,.0f} kWh'],
        ['Peak Month', f'{total.idxmax()} ({total.max():,.0f} kWh)'],
        ['Min Month', f'{total.idxmin()} ({total.min():,.0f} kWh)'],
        ['Efficiency Avg', f'{efficiency.mean():.1f}%'],
        ['Cost Estimate', f'EUR{energy_data.sum().sum() * 0.15:,.0f}']
    ]
    
    table = ax7.table(cellText=summary_data, cellLoc='left',
                     loc='center', colWidths=[0.5, 0.5])
    table.auto_set_font_size(False)
    table.set_fontsize(9)
    table.scale(1, 1.5)
    
    # Style header row
    for i in range(2):
        table[(0, i)].set_facecolor('#40466e')
        table[(0, i)].set_text_props(weight='bold', color='white')
    
    # Alternate row colors
    for i in range(1, len(summary_data)):
        for j in range(2):
            if i % 2 == 0:
                table[(i, j)].set_facecolor('#f0f0f0')
    
    ax7.set_title('Summary Statistics', fontsize=11, fontweight='bold', pad=20)
    
    plt.show()
    
    print("Complex dashboard created with:")
    print("- GridSpec for precise layout control")
    print("- 7 different visualization types")
    print("- Consistent styling across all plots")
    print("- Summary statistics table")
    print("- Interactive elements (hover values)")
    
    # Outcome:
    # Complex dashboard created with:
    # - GridSpec for precise layout control
    # - 7 different visualization types
    # - Consistent styling across all plots
    # - Summary statistics table
    # - Interactive elements (hover values)
    # [Complex multi-panel dashboard appears]
}{Creating a professional dashboard with GridSpec}

% ====================================================================
% PART VI – CUSTOM STYLES AND THEMES
% ====================================================================

\hh{Creating custom visual themes}

Consistent styling across all your visualizations creates professional reports.

\code{c16_custom_styles}{python}{
    import pandas as pd
    import matplotlib.pyplot as plt
    import matplotlib as mpl
    import numpy as np
    
    # Define custom style parameters
    def set_architecture_style():
        """Custom style for architectural visualizations"""
        plt.style.use('seaborn-v0_8-darkgrid')  # Base style
        
        # Customize further
        mpl.rcParams['figure.figsize'] = (10, 6)
        mpl.rcParams['figure.dpi'] = 100
        mpl.rcParams['savefig.dpi'] = 300
        
        # Colors
        mpl.rcParams['axes.prop_cycle'] = plt.cycler(color=[
            '#2E86AB',  # Blue
            '#A23B72',  # Purple
            '#F18F01',  # Orange
            '#C73E1D',  # Red
            '#6A994E',  # Green
            '#BC4B51'   # Rose
        ])
        
        # Fonts
        mpl.rcParams['font.family'] = 'sans-serif'
        mpl.rcParams['font.sans-serif'] = ['Arial', 'Helvetica']
        mpl.rcParams['font.size'] = 10
        mpl.rcParams['axes.titlesize'] = 12
        mpl.rcParams['axes.titleweight'] = 'bold'
        mpl.rcParams['axes.labelsize'] = 10
        mpl.rcParams['legend.fontsize'] = 9
        
        # Grid
        mpl.rcParams['grid.alpha'] = 0.3
        mpl.rcParams['grid.linestyle'] = '--'
        
        # Spines
        mpl.rcParams['axes.spines.top'] = False
        mpl.rcParams['axes.spines.right'] = False
        
        # Figure
        mpl.rcParams['figure.autolayout'] = True
        mpl.rcParams['figure.facecolor'] = 'white'
    
    # Apply custom style
    set_architecture_style()
    
    # Create sample data
    df = pd.DataFrame({
        'Quarter': ['Q1', 'Q2', 'Q3', 'Q4'],
        'Residential': [320, 350, 380, 360],
        'Commercial': [280, 290, 310, 300],
        'Industrial': [150, 160, 170, 165],
        'Public': [200, 210, 220, 215]
    })
    df.set_index('Quarter', inplace=True)
    
    # Create figure with custom style
    fig, axes = plt.subplots(2, 2, figsize=(14, 10))
    fig.suptitle('Architectural Projects Analysis - Custom Theme', 
                fontsize=16, fontweight='bold')
    
    # 1. Styled bar chart
    ax1 = axes[0, 0]
    df.plot(kind='bar', ax=ax1, width=0.8)
    ax1.set_title('Project Distribution by Type')
    ax1.set_ylabel('Area (1000 m^2)')
    ax1.set_xticklabels(ax1.get_xticklabels(), rotation=0)
    ax1.legend(title='Project Type', framealpha=0.95)
    
    # 2. Styled line chart
    ax2 = axes[0, 1]
    df.plot(kind='line', ax=ax2, marker='o', linewidth=2.5)
    ax2.set_title('Quarterly Trends')
    ax2.set_ylabel('Area (1000 m^2)')
    ax2.legend(title='Project Type', loc='best')
    
    # Add shaded region
    ax2.axhspan(300, 350, alpha=0.2, color='green', label='Target Range')
    
    # 3. Styled area chart
    ax3 = axes[1, 0]
    df.plot(kind='area', ax=ax3, alpha=0.7, stacked=True)
    ax3.set_title('Cumulative Project Area')
    ax3.set_ylabel('Area (1000 m^2)')
    ax3.legend(title='Project Type', loc='upper left')
    
    # 4. Custom scatter with annotations
    ax4 = axes[1, 1]
    
    # Create scatter data
    np.random.seed(42)
    x = np.random.uniform(100, 500, 20)
    y = x * 0.8 + np.random.uniform(-50, 50, 20)
    colors = np.random.choice(range(4), 20)
    
    scatter = ax4.scatter(x, y, c=colors, s=100, alpha=0.6, 
                         cmap='tab10', edgecolors='black', linewidth=1)
    ax4.set_title('Project Size vs Budget')
    ax4.set_xlabel('Size (m^2)')
    ax4.set_ylabel('Budget (kEUR)')
    
    # Add trend line
    z = np.polyfit(x, y, 1)
    p = np.poly1d(z)
    ax4.plot(np.sort(x), p(np.sort(x)), "r--", alpha=0.8, linewidth=2)
    
    # Add custom annotation
    ax4.annotate('Optimal\nRange', xy=(300, 250), xytext=(350, 150),
                arrowprops=dict(arrowstyle='->', connectionstyle='arc3,rad=0.3'),
                bbox=dict(boxstyle='round,pad=0.3', facecolor='yellow', alpha=0.5))
    
    plt.tight_layout()
    plt.show()
    
    print("Custom theme applied with:")
    print("- Architectural color palette")
    print("- Professional fonts and sizes")
    print("- Consistent grid and spine styles")
    print("- Automatic layout adjustment")
    
    # Reset to default style
    plt.style.use('default')
    
    # Outcome:
    # Custom theme applied with:
    # - Architectural color palette
    # - Professional fonts and sizes
    # - Consistent grid and spine styles
    # - Automatic layout adjustment
    # [Four styled charts with consistent theme]
}{Creating and applying custom visual themes}

% ====================================================================
% PART VII – SAVING AND EXPORTING
% ====================================================================

\hh{Professional export for different media}

Different outputs require different export settings---screen, print, and web all have unique requirements.

\code{c16_export_professional}{python}{
    import pandas as pd
    import matplotlib.pyplot as plt
    import numpy as np
    
    # Create a professional figure
    fig, (ax1, ax2) = plt.subplots(1, 2, figsize=(12, 5))
    
    # Sample data
    data = pd.Series([45, 38, 42, 35, 48, 52],
                     index=['Zone A', 'Zone B', 'Zone C', 'Zone D', 'Zone E', 'Zone F'])
    
    # Left plot
    data.plot(kind='bar', ax=ax1, color='steelblue')
    ax1.set_title('Energy Consumption by Zone', fontsize=14, fontweight='bold')
    ax1.set_ylabel('Energy (MWh)', fontsize=12)
    ax1.set_xlabel('Zone', fontsize=12)
    ax1.grid(True, alpha=0.3)
    
    # Right plot
    colors = plt.cm.Set3(np.arange(len(data)))
    data.plot(kind='pie', ax=ax2, autopct='%1.1f%%', colors=colors)
    ax2.set_title('Energy Distribution', fontsize=14, fontweight='bold')
    
    plt.suptitle('Building Energy Analysis 2024', fontsize=16, fontweight='bold')
    plt.tight_layout()
    
    # EXPORT OPTIONS
    
    # 1. High-resolution PNG for print (300+ DPI)
    fig.savefig('energy_print.png', 
               dpi=300,                    # High resolution
               bbox_inches='tight',        # Remove whitespace
               facecolor='white',          # White background
               edgecolor='none')           # No border
    print("Saved: energy_print.png (300 DPI for print)")
    
    # 2. Web-optimized PNG (72-150 DPI, smaller file)
    fig.savefig('energy_web.png',
               dpi=100,                    # Web resolution
               bbox_inches='tight',
               facecolor='white',
               optimize=True)              # Optimize file size
    print("Saved: energy_web.png (100 DPI for web)")
    
    # 3. Vector PDF for publications
    fig.savefig('energy_publication.pdf',
               format='pdf',               # Vector format
               bbox_inches='tight',
               facecolor='white',
               edgecolor='none')
    print("Saved: energy_publication.pdf (vector for publications)")
    
    # 4. SVG for web and editing
    fig.savefig('energy_editable.svg',
               format='svg',               # Scalable vector
               bbox_inches='tight',
               facecolor='white')
    print("Saved: energy_editable.svg (editable vector)")
    
    # 5. Transparent background for presentations
    fig.savefig('energy_transparent.png',
               dpi=150,
               bbox_inches='tight',
               transparent=True,           # Transparent background
               facecolor='none')
    print("Saved: energy_transparent.png (transparent for slides)")
    
    # 6. EPS for LaTeX documents
    fig.savefig('energy_latex.eps',
               format='eps',               # Encapsulated PostScript
               bbox_inches='tight')
    print("Saved: energy_latex.eps (for LaTeX documents)")
    
    # Show file sizes
    import os
    print("\nFile sizes:")
    for filename in ['energy_print.png', 'energy_web.png', 
                     'energy_publication.pdf', 'energy_editable.svg']:
        if os.path.exists(filename):
            size = os.path.getsize(filename) / 1024  # Convert to KB
            print(f"  {filename}: {size:.1f} KB")
    
    plt.show()
    
    print("\nExport guidelines:")
    print("- Print: 300+ DPI PNG or PDF")
    print("- Web: 72-150 DPI PNG or SVG")
    print("- Publications: PDF or EPS (vector)")
    print("- Presentations: PNG with transparency")
    print("- Further editing: SVG format")
    
    # Outcome:
    # Saved: energy_print.png (300 DPI for print)
    # Saved: energy_web.png (100 DPI for web)
    # Saved: energy_publication.pdf (vector for publications)
    # Saved: energy_editable.svg (editable vector)
    # Saved: energy_transparent.png (transparent for slides)
    # Saved: energy_latex.eps (for LaTeX documents)
    # 
    # File sizes:
    #   energy_print.png: 45.2 KB
    #   energy_web.png: 18.3 KB
    #   energy_publication.pdf: 12.7 KB
    #   energy_editable.svg: 28.9 KB
    # 
    # Export guidelines:
    # - Print: 300+ DPI PNG or PDF
    # - Web: 72-150 DPI PNG or SVG
    # - Publications: PDF or EPS (vector)
    # - Presentations: PNG with transparency
    # - Further editing: SVG format
}{Professional export settings for different media}

% ====================================================================
% PART VIII – COMPLETE PROJECT EXAMPLE
% ====================================================================

\hh{Complete project: Comprehensive building analysis report}

Let's combine everything to create a professional, publication-ready analysis.

\code{c16_complete_project}{python}{
    import pandas as pd
    import matplotlib.pyplot as plt
    import numpy as np
    from datetime import datetime
    
    # Set professional style
    plt.style.use('seaborn-v0_8-whitegrid')
    plt.rcParams['font.size'] = 10
    plt.rcParams['axes.titleweight'] = 'bold'
    
    print("=" * 60)
    print("COMPREHENSIVE BUILDING ANALYSIS REPORT")
    print("=" * 60)
    
    # Generate comprehensive dataset
    np.random.seed(42)
    
    # Building specifications
    floors = ['B1', 'B2', 'G', '1', '2', '3', '4', '5', 'R']
    floor_areas = [1200, 1200, 1500, 1400, 1400, 1400, 1400, 1400, 800]
    
    # Room type distribution per floor
    room_distribution = pd.DataFrame({
        'Floor': floors,
        'Office': [0, 100, 200, 400, 450, 450, 450, 400, 0],
        'Meeting': [0, 50, 100, 150, 150, 150, 150, 150, 0],
        'Public': [200, 150, 600, 200, 150, 150, 150, 200, 200],
        'Service': [300, 300, 200, 150, 150, 150, 150, 150, 100],
        'Circulation': [700, 600, 400, 500, 500, 500, 500, 500, 500]
    })
    room_distribution.set_index('Floor', inplace=True)
    
    # Monthly performance data
    months = pd.date_range('2024-01', '2024-12', freq='M')
    performance = pd.DataFrame({
        'Energy_MWh': [45, 42, 38, 32, 28, 25, 26, 27, 30, 35, 40, 44],
        'Water_m3': [850, 820, 890, 920, 980, 1050, 1100, 1080, 950, 900, 860, 830],
        'Cost_kEUR': [62, 58, 54, 48, 44, 42, 43, 44, 47, 52, 57, 61],
        'Occupancy_pct': [78, 82, 85, 88, 90, 86, 75, 72, 88, 91, 87, 80],
        'Satisfaction': [7.2, 7.4, 7.6, 7.8, 8.0, 8.1, 7.9, 7.8, 7.7, 7.5, 7.3, 7.2]
    }, index=months)
    
    # Create comprehensive figure
    fig = plt.figure(figsize=(20, 12))
    
    # Main title with building info
    fig.text(0.5, 0.98, 'BUILDING PERFORMANCE ANALYSIS REPORT 2024',
            ha='center', fontsize=18, fontweight='bold')
    fig.text(0.5, 0.96, f'Total Area: {sum(floor_areas):,} m2 | Floors: {len(floors)} | Generated: {datetime.now().strftime("%Y-%m-%d")}',
            ha='center', fontsize=10, style='italic')
    
    # Create subplot mosaic
    gs = fig.add_gridspec(3, 4, hspace=0.3, wspace=0.25, 
                         left=0.05, right=0.95, top=0.92, bottom=0.05)
    
    # 1. Space allocation stacked bar
    ax1 = fig.add_subplot(gs[0, :2])
    room_distribution.plot(kind='bar', stacked=True, ax=ax1,
                           colormap='tab10', width=0.8)
    ax1.set_title('Space Allocation by Floor', fontsize=12)
    ax1.set_ylabel('Area (m2)', fontsize=10)
    ax1.set_xlabel('Floor', fontsize=10)
    ax1.legend(title='Space Type', bbox_to_anchor=(1.15, 1), fontsize=8)
    ax1.set_xticklabels(ax1.get_xticklabels(), rotation=0)
    
    # Add total labels
    totals = room_distribution.sum(axis=1)
    for i, (idx, total) in enumerate(totals.items()):
        ax1.text(i, total + 20, f'{total:,.0f}', ha='center', fontsize=8)
    
    # 2. Energy consumption trend with forecast
    ax2 = fig.add_subplot(gs[0, 2:])
    
    # Plot actual data
    ax2.plot(performance.index, performance['Energy_MWh'], 
            marker='o', linewidth=2, label='Actual', color='#2E86AB')
    
    # Add trend line
    x_numeric = np.arange(len(performance))
    z = np.polyfit(x_numeric, performance['Energy_MWh'], 2)
    p = np.poly1d(z)
    ax2.plot(performance.index, p(x_numeric), '--', 
            alpha=0.7, label='Trend', color='red')
    
    # Shade seasons
    ax2.axvspan(performance.index[5], performance.index[7], 
               alpha=0.1, color='yellow', label='Summer')
    ax2.axvspan(performance.index[11], performance.index[11], 
               alpha=0.1, color='lightblue')
    ax2.axvspan(performance.index[0], performance.index[1], 
               alpha=0.1, color='lightblue', label='Winter')
    
    ax2.set_title('Energy Consumption Trend', fontsize=12)
    ax2.set_ylabel('Energy (MWh)', fontsize=10)
    ax2.legend(loc='upper right', fontsize=8)
    ax2.grid(True, alpha=0.3)
    
    # Format x-axis
    import matplotlib.dates as mdates
    ax2.xaxis.set_major_formatter(mdates.DateFormatter('%b'))
    
    # 3. Cost breakdown pie chart
    ax3 = fig.add_subplot(gs[1, 0])
    
    cost_breakdown = pd.Series({
        'Energy': 35,
        'Water': 8,
        'Maintenance': 22,
        'Staff': 25,
        'Other': 10
    })
    
    colors = ['#FF9999', '#66B2FF', '#99FF99', '#FFCC99', '#FF99CC']
    wedges, texts, autotexts = ax3.pie(cost_breakdown, 
                                        labels=cost_breakdown.index,
                                        autopct='%1.1f%%',
                                        colors=colors,
                                        startangle=90,
                                        explode=[0.05, 0, 0, 0, 0])
    
    ax3.set_title('Annual Cost Distribution', fontsize=12)
    for autotext in autotexts:
        autotext.set_fontsize(8)
        autotext.set_fontweight('bold')
        autotext.set_color('white')
    
    # 4. Occupancy vs Satisfaction scatter
    ax4 = fig.add_subplot(gs[1, 1])
    
    scatter = ax4.scatter(performance['Occupancy_pct'], 
                         performance['Satisfaction'],
                         c=range(len(performance)),
                         cmap='viridis',
                         s=100, alpha=0.6,
                         edgecolors='black', linewidth=1)
    
    # Add trend line
    z = np.polyfit(performance['Occupancy_pct'], performance['Satisfaction'], 1)
    p = np.poly1d(z)
    ax4.plot(np.sort(performance['Occupancy_pct']), 
            p(np.sort(performance['Occupancy_pct'])), 
            'r--', alpha=0.7)
    
    ax4.set_title('Occupancy vs Satisfaction', fontsize=12)
    ax4.set_xlabel('Occupancy (%)', fontsize=10)
    ax4.set_ylabel('Satisfaction Score', fontsize=10)
    ax4.grid(True, alpha=0.3)
    
    # Add colorbar for months
    cbar = plt.colorbar(scatter, ax=ax4)
    cbar.set_label('Month', fontsize=8)
    cbar.ax.tick_params(labelsize=8)
    
    # 5. Monthly comparison heatmap
    ax5 = fig.add_subplot(gs[1, 2:])
    
    # Normalize data for heatmap
    perf_normalized = performance[['Energy_MWh', 'Water_m3', 'Cost_kEUR', 'Occupancy_pct']].T
    perf_normalized = (perf_normalized - perf_normalized.mean(axis=1).values.reshape(-1, 1)) / perf_normalized.std(axis=1).values.reshape(-1, 1)
    
    im = ax5.imshow(perf_normalized, cmap='RdBu_r', aspect='auto', vmin=-2, vmax=2)
    ax5.set_title('Monthly Performance Heatmap (Normalized)', fontsize=12)
    ax5.set_xticks(range(12))
    ax5.set_xticklabels(['J', 'F', 'M', 'A', 'M', 'J', 'J', 'A', 'S', 'O', 'N', 'D'])
    ax5.set_yticks(range(4))
    ax5.set_yticklabels(['Energy', 'Water', 'Cost', 'Occupancy'])
    
    # Add colorbar
    cbar = plt.colorbar(im, ax=ax5)
    cbar.set_label('Std Dev from Mean', fontsize=8)
    cbar.ax.tick_params(labelsize=8)
    
    # 6. Box plot comparison
    ax6 = fig.add_subplot(gs[2, :2])
    
    # Create seasonal data
    seasonal_energy = pd.DataFrame({
        'Winter': performance['Energy_MWh'].iloc[[0, 1, 11]],
        'Spring': performance['Energy_MWh'].iloc[[2, 3, 4]],
        'Summer': performance['Energy_MWh'].iloc[[5, 6, 7]],
        'Fall': performance['Energy_MWh'].iloc[[8, 9, 10]]
    })
    
    bp = ax6.boxplot([seasonal_energy[col].values for col in seasonal_energy.columns],
                     labels=seasonal_energy.columns,
                     patch_artist=True)
    
    # Color the boxes
    colors = ['lightblue', 'lightgreen', 'yellow', 'orange']
    for patch, color in zip(bp['boxes'], colors):
        patch.set_facecolor(color)
        patch.set_alpha(0.7)
    
    ax6.set_title('Seasonal Energy Distribution', fontsize=12)
    ax6.set_ylabel('Energy (MWh)', fontsize=10)
    ax6.grid(True, axis='y', alpha=0.3)
    
    # 7. Summary metrics
    ax7 = fig.add_subplot(gs[2, 2:])
    ax7.axis('off')
    
    # Calculate KPIs
    kpis = [
        ['KEY PERFORMANCE INDICATORS', '2024'],
        ['', ''],
        ['Total Energy Consumption', f'{performance["Energy_MWh"].sum():.0f} MWh'],
        ['Average Monthly Energy', f'{performance["Energy_MWh"].mean():.1f} MWh'],
        ['Total Water Usage', f'{performance["Water_m3"].sum():,.0f} m3'],
        ['Total Annual Cost', f'EUR{performance["Cost_kEUR"].sum():.0f}k'],
        ['Average Occupancy', f'{performance["Occupancy_pct"].mean():.1f}%'],
        ['Average Satisfaction', f'{performance["Satisfaction"].mean():.1f}/10'],
        ['Energy per m2', f'{performance["Energy_MWh"].sum()*1000/sum(floor_areas):.1f} kWh/m2'],
        ['Cost per m2', f'EUR{performance["Cost_kEUR"].sum()*1000/sum(floor_areas):.2f}/m2']
    ]
    
    table = ax7.table(cellText=kpis, cellLoc='left',
                     loc='center', colWidths=[0.6, 0.4])
    table.auto_set_font_size(False)
    table.set_fontsize(9)
    table.scale(1, 1.8)
    
    # Style the table
    for i in range(len(kpis)):
        if i == 0:  # Header row
            for j in range(2):
                table[(i, j)].set_facecolor('#2E86AB')
                table[(i, j)].set_text_props(weight='bold', color='white')
        elif i % 2 == 0:  # Alternate rows
            for j in range(2):
                table[(i, j)].set_facecolor('#f0f0f0')
    
    ax7.set_title('Performance Summary', fontsize=12, pad=20)
    
    # Save the complete figure
    plt.savefig('building_analysis_report.pdf', 
               dpi=300, bbox_inches='tight', facecolor='white')
    plt.savefig('building_analysis_report.png', 
               dpi=150, bbox_inches='tight', facecolor='white')
    
    plt.show()
    
    print("\n" + "=" * 60)
    print("REPORT SUMMARY")
    print("=" * 60)
    
    # Print summary statistics
    print(f"\nBuilding Overview:")
    print(f"  Total Floors: {len(floors)}")
    print(f"  Total Area: {sum(floor_areas):,} m2")
    print(f"  Average Floor Area: {np.mean(floor_areas):,.0f} m2")
    
    print(f"\nSpace Allocation:")
    total_by_type = room_distribution.sum()
    for space_type, area in total_by_type.items():
        pct = area / total_by_type.sum() * 100
        print(f"  {space_type:12s}: {area:5.0f} m2 ({pct:4.1f}%)")
    
    print(f"\nAnnual Performance:")
    print(f"  Energy: {performance['Energy_MWh'].sum():.0f} MWh")
    print(f"  Water: {performance['Water_m3'].sum():,.0f} m3")
    print(f"  Cost: EUR{performance['Cost_kEUR'].sum():.0f}k")
    print(f"  Avg Occupancy: {performance['Occupancy_pct'].mean():.1f}%")
    print(f"  Avg Satisfaction: {performance['Satisfaction'].mean():.1f}/10")
    
    print(f"\nEfficiency Metrics:")
    print(f"  Energy Intensity: {performance['Energy_MWh'].sum()*1000/sum(floor_areas):.1f} kWh/m2/year")
    print(f"  Cost per Area: EUR{performance['Cost_kEUR'].sum()*1000/sum(floor_areas):.2f}/m2/year")
    print(f"  Water per Area: {performance['Water_m3'].sum()/sum(floor_areas):.1f} m3/m2/year")
    
    print(f"\nFiles Generated:")
    print(f"  - building_analysis_report.pdf (print quality)")
    print(f"  - building_analysis_report.png (digital display)")
    
    # Outcome:
    # ============================================================
    # COMPREHENSIVE BUILDING ANALYSIS REPORT
    # ============================================================
    # 
    # ============================================================
    # REPORT SUMMARY
    # ============================================================
    # 
    # Building Overview:
    #   Total Floors: 9
    #   Total Area: 11,100 m2
    #   Average Floor Area: 1,233 m2
    # 
    # Space Allocation:
    #   Office      : 2450 m2 (22.1%)
    #   Meeting     :  900 m2 ( 8.1%)
    #   Public      : 2000 m2 (18.0%)
    #   Service     : 1450 m2 (13.1%)
    #   Circulation : 4300 m2 (38.7%)
    # 
    # Annual Performance:
    #   Energy: 432 MWh
    #   Water: 10,880 m3
    #   Cost: EUR612k
    #   Avg Occupancy: 83.5%
    #   Avg Satisfaction: 7.6/10
    # 
    # Efficiency Metrics:
    #   Energy Intensity: 38.9 kWh/m2/year
    #   Cost per Area: EUR55.14/m2/year
    #   Water per Area: 1.0 m3/m2/year
    # 
    # Files Generated:
    #   - building_analysis_report.pdf (print quality)
    #   - building_analysis_report.png (digital display)
    # [Comprehensive multi-panel report visualization]
}{Complete building analysis report with professional visualizations}